\documentclass[../distribution_theory_notes.tex]{subfiles}
\begin{document}
\section{Aula 20 - 04 de Novembro, 2024}
\subsection{Motivações}
\begin{itemize}
	\item Distribuições com suporte concentrado em um ponto;
	\item Densidade das funções teste em \(\mathcal{D}'(\Omega )\) para um aberto arbitrário \(\Omega \subseteq \mathbb{R}^{n}\);
	\item Preparativos para o produto de convolução de duas distribuições.
\end{itemize}

\subsection{Distribuições com Suporte Compacto e Distribuições Suportadas em um Ponto}
O objetivo da aula passada foi identificar o espaço \(\mathcal{E}'(\Omega )\), dual topológico de \(\mathcal{C}^{\infty}(\Omega )\), sendo este último munido da topologia da convergência uniforme em compactos de \(\Omega \) das funções e de suas derivadas, com um subespaço de \(\mathcal{D}'(\Omega )\): precisamente, aquele formado pelas distribuições que possuem suporte compacto. Além disso, aprendemos como usar partições da unidade para ``reconstruir'' uma distribuição a partir de suas localizações.

A fim de estabelecermos a identificação acima, por um lado, dada \(u\in \mathcal{E}'(\Omega )\), existem \(K\Subset \Omega \), \(c>0\) e \(m\in \mathbb{Z}_{+}\) tais que
\[
	|\langle u, f \rangle |\leq c\sum\limits_{|\alpha |\leq m}^{}\sup_{x\in K}|\partial^{\alpha }f(x)|,\quad \forall f\in \mathcal{C}^{\infty}(\Omega ). \qquad (\#)
\]
Com efeito, a família de seminormas
\[
	p_{(m,K)}(f)=\sum\limits_{|\alpha |\leq m}^{}\sup_{x\in K}|\partial^{\alpha }f(x)|,
\]
para \(K\subset \Omega \) e \(m\in \mathbb{Z}_{+}\), define a topologia de \(\mathcal{C}^{\infty}(\Omega )\). Da estimativa \((\#)\), se \(\varphi_{j}\to 0\) em \(\mathcal{C}_{c}^{\infty}(\Omega )\), então
\[
	p_{(m,K)}(\varphi_{j})\to 0 \Rightarrow \langle u, \varphi_{j} \rangle\to 0,
\]
mostrando que a restrição \(v\coloneqq u|_{\mathcal{C}_{c}^{\infty}(\Omega )}\) pertence a \(\mathcal{D}'(\Omega )\). Além disso, se \(\varphi \in \mathcal{C}_{c}^{\infty}(\Omega )\) satisfaz \(\mathrm{supp}(\varphi )\cap K=\emptyset \), então \(p_{(m,K)}(\varphi )=0\), donde
\[
	\langle v, \varphi  \rangle=0,
\]
e, consequentemente,
\[
	\mathrm{supp}(v)\subseteq K,
\]
de modo que \(v\) possui suporte compacto contido em \(\Omega \).

Por outro lado, dada \(u\in \mathcal{D}'(\Omega )\) com suporte compacto K, \(\mathrm{supp}(u)=K\subset \Omega \), fixemos uma função de Urysohn \(\psi \in \mathcal{C}_{c}^{\infty}(\Omega )\) associada a K, isto é, \(\psi \equiv 1\) em uma vizinhança de K, e definamos
\begin{align*}
	\overline{u}: & \mathcal{C}^{\infty}(\Omega )\longrightarrow \mathbb{C}                         \\
	              & f\longmapsto \langle \overline{u}, f\rangle \coloneqq \langle u, \psi f\rangle.
\end{align*}
Vimos que \(\overline{u}\) está bem definida independentemente da escolha de \(\psi \), que \(\overline{u}\in \mathcal{E}'(\Omega )\) e que ela é a única extensão de u a \(\mathcal{C}^{\infty}(\Omega )\) pertencente a \(\mathcal{E}'(\Omega )\), pois \(\mathcal{C}_{c}^{\infty}(\Omega )\) é denso em \(\mathcal{C}^{\infty}(\Omega )\); mais precisamente, se \(K_{0}=\mathrm{supp}(\psi )\), então existem \(c>0\) e \(m\in \mathbb{Z}_{+}\) tais que
\[
	|\langle \overline{u}, f\rangle |\leq c\sum\limits_{|\alpha |\leq m}^{}\sup_{x\in K_{0}}|\partial^{\alpha }f(x)|,\quad \forall f\in \mathcal{C}^{\infty}(\Omega ),
\]
confirmando, mais uma vez, que \(\overline{u}\in \mathcal{E}'(\Omega )\).

Queremos agora utilizar parte das ideias acima para estabelecer um dos fatos mais relevantes da teoria, segundo o qual qualquer distribuição com suporte sendo \(\{0\}\)\footnote{Basta ser um conjunto unitário, ou tem que ser suportada no 0? Talvez um argumento de translação responda isso.} pode ser escrita como uma combinação linear de um número finito de derivadas da distribuição \(\delta \).

Dependendo do tempo, ampliaremos ainda para um aberto arbitrário \(\Omega \subseteq \mathbb{R}^{n}\) o resultado de densidade das funções teste em \(\mathcal{D}'(\mathbb{R}^{n})\), demonstrado anteriormente, e iniciaremos os preparativos para a terceira parte da convolução, em que ambos os fatores são distribuições.

Sabemos que \(\mathrm{supp}(\delta )=\{0\}\) e, como
\[
	\mathrm{supp}(\partial^{\alpha }u)\subseteq \mathrm{supp}(u),\quad \forall u\in \mathcal{D}'(\Omega ),\; \alpha \in \mathbb{Z}_{+}^{n},
\]
temos, em particular,
\[
	\mathrm{supp}(\partial^{\alpha }\delta )\subseteq \{0\},\quad \forall \alpha \in \mathbb{Z}_{+}^{n}.
\]
Considerando os monômios
\[
	\psi_{\alpha }(x)=(-1)^{|\alpha |}\frac{x^{\alpha }}{\alpha !},
\]
(ou versões deles multiplicadas por uma função de Urysohn ao redor da origem), veremos que
\[
	\langle \partial ^{\alpha }\delta , \psi_{\alpha } \rangle=1,
\]
e portanto, para qualquer mmjk, \(\partial^{\alpha }\delta \neq 0\), levando à conclusão de que o suporte de \(\partial^{\alpha }\delta \) não pode ser vazio, donde segue que ele necessariamente deve ser sempre o conjunto \(\{0\}\).

Partindo destas funções \(\psi_{\alpha }\), vamos mostrar que, mais geralmente, se \(u = \sum\limits_{| \alpha  |\leq m}^{}a_{\alpha }\partial^{\alpha }\delta \) tem algum \(a_{\alpha }\neq 0\), então o suporte de \(u\) deve ser \(\{0\}\)
\begin{tcolorbox}[
		skin=enhanced,
		title=Observação,
		fonttitle=\bfseries,
		colframe=black,
		colbacktitle=cyan!75!white,
		colback=cyan!15,
		colbacklower=black,
		coltitle=black,
		drop fuzzy shadow,
		%drop large lifted shadow
	]
	Para as funções \(\psi_{\alpha }(x)=(-1)^{|\alpha |}x^{\alpha }/\alpha !\), as derivadas de \(\delta \) satisfazem um ``delta de Kronecker'':
	\[
		\langle \partial^{\beta }\delta , \psi_{\alpha }\rangle
		=(-1)^{|\beta |}\partial^{\beta }\psi_{\alpha }(0)
		=\begin{cases}
			1, & \beta =\alpha ,     \\
			0, & \beta \neq \alpha .
		\end{cases}
	\]
\end{tcolorbox}

\begin{theorem*}
	Seja \(u\in \mathcal{D}'(\mathbb{R}^{n})\) com \(\mathrm{supp}(u)=\{0\}\). Existem \(m\in \mathbb{Z}_{+}\) e escalares \(a_{\alpha }\in \mathbb{C}\), com \(|\alpha |\leq m\), tais que
	\[
		u=\sum\limits_{|\alpha |\leq m}^{}a_{\alpha }\partial^{\alpha }\delta .
	\]
\end{theorem*}

A demonstração possui alguns pontos-chave, um deles sendo a seguinte forma da Fórmula de Taylor.
\begin{theorem*}[Fórmula de Taylor com Resto Integral]
	Se \(f\in \mathcal{C}^{m+1}(\mathbb{R}^{n})\), com \(m=0,1,2,\dotsc \), então, para todo \(x\in \mathbb{R}^{n}\),
	\[
		f(x)=\sum\limits_{|\alpha |\leq m}^{}\frac{\partial^{\alpha }f(0)}{\alpha !}x^{\alpha }+R_{m}(x),
	\]
	em que
	\[
		R_{m}(x)=(m+1)\sum\limits_{|\alpha |=m+1}^{}\int_{0}^{1}(1-t)^{m}\frac{\partial^{\alpha }f(tx)}{\alpha !}x^{\alpha }\mathrm{d}t
	\]
	é o resto integral.
\end{theorem*}
O fato crucial para o que segue é que, quando \(f\in \mathcal{C}^{\infty}(\mathbb{R}^{n})\), o resto \(R_{m}\) também pertence a \(\mathcal{C}^{\infty}(\mathbb{R}^{n})\), sendo suficiente uma derivação sob o sinal de integral pela Regra de Leibniz.

\begin{proof*}
	Seja \(f\in \mathcal{C}^{\infty}(\mathbb{R}^{n})\) qualquer; pelo teorema anterior, para todo \(N\in \mathbb{N}\), podemos escrever
	\[
		f(x)=\sum\limits_{|\alpha |\leq N}^{}\frac{\partial^{\alpha }f(0)}{\alpha !}x^{\alpha }+R_{N}(x). \qquad (I)
	\]
	Como \(\mathrm{supp}(u)=\{0\}\), identificamos u com sua extensão a \(\mathcal{E}'(\mathbb{R}^{n})\). Se \(\psi \in \mathcal{C}_{c}^{\infty}(\mathbb{R}^{n})\) é qualquer função de Urysohn associada a \(\{0\}\), então
	\[
		\langle u, f\rangle=\langle u, \psi f\rangle.
	\]
	Decompondo \(f\) de acordo com \((I)\), obtemos
	\begin{align*}
		\langle u, f\rangle
		 & =\biggl\langle u, \sum\limits_{|\alpha |\leq N}^{}\frac{\partial^{\alpha }f(0)}{\alpha !}\psi x^{\alpha }+\psi R_{N}\biggr\rangle        \\
		 & =\sum\limits_{|\alpha |\leq N}^{}\frac{\partial^{\alpha }f(0)}{\alpha !}\langle u, \psi x^{\alpha }\rangle+\langle u, \psi R_{N}\rangle.
	\end{align*}
	Como
	\[
		\langle \partial^{\alpha }\delta , f\rangle=(-1)^{|\alpha |}\partial^{\alpha }f(0)
	\]
	e \(\langle u, \psi x^{\alpha }\rangle\) independe da escolha da função de Urysohn \(\psi \), podemos definir
	\[
		a_{\alpha }\coloneqq \frac{(-1)^{|\alpha |}}{\alpha !}\langle u, \psi x^{\alpha }\rangle,
	\]
	para \(\alpha \in \mathbb{Z}_{+}^{n}\). Assim,
	\[
		\langle u, f\rangle
		=\sum\limits_{|\alpha |\leq N}^{}a_{\alpha }\langle \partial^{\alpha }\delta , f\rangle+\langle u, \psi R_{N}\rangle
		=\biggl\langle \sum\limits_{|\alpha |\leq N}^{}a_{\alpha }\partial^{\alpha }\delta , f\biggr\rangle+\langle u, \psi R_{N}\rangle.
	\]
	Logo, o teorema estará demonstrado se encontrarmos N e uma função de Urysohn conveniente para os quais
	\[
		\langle u, \psi R_{N}\rangle=0.
	\]

	Com efeito, fixemos \(\varphi \in \mathcal{C}_{c}^{\infty}(\overline{B}(0;1))\) tal que \(\varphi \geq 0\) e \(\int \varphi =1\), e coloquemos
	\[
		\varphi_{\varepsilon }(x)=\varepsilon^{-n}\varphi \biggl(\frac{x}{\varepsilon }\biggr).
	\]
	Para \(\varepsilon >0\), definimos
	\[
		\psi_{\varepsilon }(x)
		\coloneqq \varepsilon^{-n}\int_{\overline{B}(0;2\varepsilon )}^{}\varphi \biggl(\frac{x-y}{\varepsilon }\biggr)\mathrm{d}y
		=\varphi_{\varepsilon }*\chi_{\overline{B}(0;2\varepsilon )}(x).
	\]
	Examinando a demonstração do Lema de Urysohn, vemos que \(\psi_{\varepsilon }\) é uma função de Urysohn associada a \(K=\{0\}\); de fato,
	\[
		\mathrm{supp}(\psi_{\varepsilon })
		\subseteq \overline{B}(0;\varepsilon )+\overline{B}(0;2\varepsilon )
		=\overline{B}(0;3\varepsilon )
	\]
	e
	\[
		\psi_{\varepsilon }(x)=1,\quad \forall x\in B(0;\varepsilon ).
	\]

	Passamos às estimativas: para \(\beta \in \mathbb{Z}_{+}^{n}\),
	\begin{align*}
		\partial^{\beta }\psi_{\varepsilon }(x)
		 & =\varepsilon^{-n-|\beta |}\int_{\overline{B}(0;2\varepsilon )}^{}(\partial^{\beta }\varphi )\biggl(\frac{x-y}{\varepsilon }\biggr)\mathrm{d}y,
	\end{align*}
	donde
	\[
		\|\partial^{\beta }\psi_{\varepsilon }\|_{\infty }
		\leq 2^{n}B_{n}\varepsilon^{-|\beta |}\|\partial^{\beta }\varphi \|_{\infty },
	\]
	onde \(B_{n}=\mathrm{vol}(\overline{B}(0;1))\).

	Por outro lado, lembrando que
	\[
		\partial^{\eta }x^{\alpha }
		=\frac{\alpha !}{(\alpha -\eta )!}x^{\alpha -\eta }
	\]
	quando \(\eta_{j}\leq \alpha_{j}\) para todo \(j=1,\dotsc ,n\), e que \(\partial^{\eta }x^{\alpha }=0\) se algum \(\eta_{j}>\alpha_{j}\), a Regra de Leibniz fornece
	\begin{align*}
		\partial^{\beta }R_{N}(x)
		 & =(N+1)\sum\limits_{|\alpha |=N+1}^{}\int_{0}^{1}(1-t)^{N}\partial^{\beta }\biggl(\frac{\partial^{\alpha }f(tx)}{\alpha !}x^{\alpha }\biggr)\mathrm{d}t \\
		 & =(N+1)\sum\limits_{|\alpha |=N+1}^{}\int_{0}^{1}(1-t)^{N}
		\sum\limits_{\substack{\gamma +\eta =\beta                                                                                                                \\ \eta \leq \alpha }}^{}
		\frac{\beta !}{\gamma !\eta !}\frac{t^{|\gamma |}\partial^{\gamma +\alpha }f(tx)}{\alpha !}
		\frac{\alpha !}{(\alpha -\eta )!}x^{\alpha -\eta }\mathrm{d}t.
	\end{align*}
	Consequentemente, para \(|x|\leq 1\) e \(|\beta |\leq N+1\), existe uma constante \(C_{N}^{(\beta )}>0\) tal que
	\[
		|\partial^{\beta }R_{N}(x)|\leq C_{N}^{(\beta )}|x|^{N+1-|\beta |}. \qquad (II)
	\]
	Aqui, as constantes dependem apenas de N, \(\beta \) e de uma quantidade finita de derivadas de f sobre \(\overline{B}(0;1)\).
	O destrinchar das contas, levemente assustador, é:
	\begin{align*}
		 & | \partial^{\beta }R_{N}(x) | \leq (N+1)\sum\limits_{| \alpha  |=N+1}^{}\int_{0}^{1}\overbrace{(1-t)^{N}}^{\mathclap{\in[0, 1]}}\biggl[\sum\limits_{\gamma +\eta =\beta }^{}t^{| \gamma  |}\frac{\beta !}{\gamma !\eta !}\frac{\alpha !}{(\alpha -\eta )!}\biggl\vert \frac{\partial^{\gamma +\alpha }f}{\alpha !}(tx) \biggr\vert | x |^{\alpha - | \eta  |}\biggr] \mathrm{dt}                                                                                             \\
		 & \leq (N+1)\sum\limits_{| \alpha  | = N+1}^{}\sum\limits_{\gamma +\eta =\beta }^{}\frac{\beta !}{\gamma !\eta !}\frac{\alpha !}{(\alpha -\eta )!}\int_{0}^{1}\biggl\vert \frac{\partial^{\gamma +\alpha }f}{\alpha !}(tx) \biggr\vert \mathrm{d}t | x |^{| \alpha  |-| \eta  |}                                                                                                                                                                                               \\
		 & \leq (N+1)\sum\limits_{| \alpha  |=N+1}^{}\biggl[\sum\limits_{\gamma +\eta =\beta }^{}\underbrace{\frac{\beta !}{\gamma !\eta !}\frac{\alpha !}{(\alpha -\eta )!}}_{\mathclap{\coloneqq C(\alpha , \beta , \gamma , \eta )}}\int_{0}^{1}\underbrace{\biggl\vert \frac{\partial^{\gamma +\alpha }f}{\alpha !}(tx) \biggr\vert}_{\mathclap{\leq 1/\alpha ! \Vert \partial^{\gamma +\alpha }f \Vert_{L^{\infty}(\overline{B}(0; 1))}}} \mathrm{d}t\biggr]| x |^{N+1-| \beta  |}
	\end{align*}

	Como \(K_{0}=\overline{B}(0;1)\) é compacto, existem \(c>0\) e \(m\in \mathbb{Z}_{+}\) tais que
	\[
		|\langle u, \phi \rangle |
		\leq c\sum\limits_{|\beta |\leq m}^{}\sup_{x\in \overline{B}(0;1)}|\partial^{\beta }\phi (x)|,
		\quad \forall \phi \in \mathcal{C}_{c}^{\infty}(\overline{B}(0;1)).
	\]
	Se \(0<\varepsilon <\frac{1}{3}\), então
	\[
		\mathrm{supp}(\psi_{\varepsilon }R_{N})\subseteq \overline{B}(0;3\varepsilon )\subseteq \overline{B}(0;1),
	\]
	e, portanto,
	\[
		|\langle u, \psi_{\varepsilon }R_{N}\rangle |
		\leq c\sum\limits_{|\beta |\leq m}^{}\sup_{|x|\leq 3\varepsilon }|\partial^{\beta }(\psi_{\varepsilon }R_{N})(x)|. \qquad (\#\#)
	\]
	Aplicando novamente a Regra de Leibniz e as estimativas anteriores, para \(|\beta |\leq N+1\) e \(|x|\leq 3\varepsilon \), obtemos
	\begin{align*}
		|\partial^{\beta }(\psi_{\varepsilon }R_{N})(x)|
		 & \leq \sum\limits_{\gamma +\eta =\beta }^{}\frac{\beta !}{\gamma !\eta !}
		|\partial^{\gamma }\psi_{\varepsilon }(x)|\,|\partial^{\eta }R_{N}(x)|      \\
		 & \leq \sum\limits_{\gamma +\eta =\beta }^{}\frac{\beta !}{\gamma !\eta !}
		\bigl[2^{n}B_{n}\|\partial^{\gamma }\varphi \|_{\infty }\varepsilon^{-|\gamma |}\bigr]
		\bigl[C_{N}^{(\eta )}\varepsilon^{N+1-|\eta |}\bigr]                        \\
		 & \leq d_{(m,N)}\varepsilon^{N+1-|\beta |},
	\end{align*}
	para alguma constante \(d_{(m,N)}>0\).

	Finalmente, como em \((\#\#)\) a ordem máxima de derivação é m, escolhemos \(N=m\). Então, para \(|\beta |\leq m\) e \(0<\varepsilon <1\),
	\[
		\varepsilon^{m+1-|\beta |}\leq \varepsilon,
	\]
	de modo que
	\[
		\sup_{|x|\leq 3\varepsilon }|\partial^{\beta }(\psi_{\varepsilon }R_{m})(x)|\leq d_{(n,m)}\varepsilon .
	\]
	Voltando a \((\#\#)\), resulta que, para alguma constante \(\widetilde{d}_{(n,m)}>0\),
	\[
		|\langle u, \psi_{\varepsilon }R_{m}\rangle |\leq \widetilde{d}_{(n,m)}\varepsilon .
	\]
	Entretanto, \(\psi_{\varepsilon }\) é função de Urysohn associada a \(\mathrm{supp}(u)=\{0\}\), logo o valor
	\[
		\langle u, \psi_{\varepsilon }R_{m}\rangle=\langle u, R_{m}\rangle
	\]
	independe de \(\varepsilon \). Fazendo \(\varepsilon \to 0^{+}\), concluímos que
	\[
		\langle u, R_{m}\rangle=0,
	\]
	e, portanto,
	\[
		u=\sum\limits_{|\alpha |\leq m}^{}a_{\alpha }\partial^{\alpha }\delta . \text{ \qedsymbol}
	\]
\end{proof*}

\subsection{Densidade das Funções Teste em \(\mathcal{D}'(\Omega )\)}
Queremos demonstrar que \(\mathcal{C}_{c}^{\infty}(\Omega )\) é denso, na topologia fraca-*, em \(\mathcal{D}'(\Omega )\), para um aberto arbitrário \(\Omega \subseteq \mathbb{R}^{n}\). Grande parte das ideias necessárias já apareceu na demonstração do mesmo resultado para \(\Omega =\mathbb{R}^{n}\); por isso, o caso geral segue essencialmente de algumas observações que permitem ajustar aquela prova.

\begin{theorem*}
	Se \(\Omega \subseteq \mathbb{R}^{n}\) é um aberto, então \(\mathcal{C}_{c}^{\infty}(\Omega )\) é denso em \(\mathcal{D}'(\Omega )\) na topologia fraca-*.
\end{theorem*}
\begin{proof*}
	Começamos com algumas observações.

	(1) Se \(f\in \mathcal{C}^{\infty}(\Omega )\) e \(u\in \mathcal{D}'(\Omega )\), então
	\[
		\mathrm{supp}(fu)\subseteq \mathrm{supp}(f)\cap \mathrm{supp}(u).
	\]
	Com efeito, se \(\varphi \in \mathcal{C}_{c}^{\infty}(\Omega \setminus{\mathrm{supp}(u)})\), então
	\[
		\mathrm{supp}(f\varphi )\subseteq \mathrm{supp}(f)\cap \mathrm{supp}(\varphi )
	\]
	e \(f\varphi \in \mathcal{C}_{c}^{\infty}(\Omega \setminus{\mathrm{supp}(u)})\). Logo,
	\[
		\langle fu, \varphi  \rangle=\langle u, f\varphi  \rangle=0,
	\]
	donde \(\mathrm{supp}(fu)\subseteq \mathrm{supp}(u)\). Analogamente, se \(\varphi \in \mathcal{C}_{c}^{\infty}(\Omega \setminus{\mathrm{supp}(f)})\), então \(f\varphi =0\), e assim \(\mathrm{supp}(fu)\subseteq \mathrm{supp}(f)\).

	(2) Seja \((K_{j})_{j}\) um esgotamento de \(\Omega \), e seja \(\psi_{j}\in \mathcal{C}_{c}^{\infty}(\Omega )\) uma função de Urysohn associada a \(K_{j}\). Se \(\varphi \in \mathcal{C}^{\infty}(\Omega )\), então
	\[
		\psi_{j}\varphi \stackrel{j\to \infty}\longrightarrow \varphi \quad \text{em }\mathcal{C}^{\infty}(\Omega ).
	\]
	Portanto, \(\mathcal{C}_{c}^{\infty}(\Omega )\) é denso em \(\mathcal{C}^{\infty}(\Omega )\).

	(3) Se \(u\in \mathcal{D}'(\Omega )\) e \(\psi \in \mathcal{C}_{c}^{\infty}(\Omega )\), então
	\[
		\psi u\in \mathcal{E}'(\Omega ),
	\]
	pois, por (1),
	\[
		\mathrm{supp}(\psi u)\subseteq \mathrm{supp}(\psi ),
	\]
	que é compacto.

	(4) O espaço \(\mathcal{E}'(\Omega )\) é denso em \(\mathcal{D}'(\Omega )\). Com efeito, se \(u\in \mathcal{D}'(\Omega )\) e \((\psi_{j})_{j}\) é como em (2), então \(\psi_{j}u\in \mathcal{E}'(\Omega )\) e, para toda \(\varphi \in \mathcal{C}_{c}^{\infty}(\Omega )\),
	\[
		\langle \psi_{j}u, \varphi  \rangle
		=\langle u, \psi_{j}\varphi  \rangle
		\stackrel{j\to \infty}\longrightarrow \langle u, \varphi  \rangle.
	\]
	Assim,
	\[
		\psi_{j}u\stackrel{j\to \infty}\longrightarrow u\quad \text{em }\mathcal{D}'(\Omega ).
	\]

	(5) Resta mostrar que \(\mathcal{C}_{c}^{\infty}(\Omega )\) é denso em \(\mathcal{E}'(\Omega )\), este último naturalmente munido da topologia fraca-* como subespaço de \(\mathcal{D}'(\Omega )\). Seja, então, \(u\in \mathcal{E}'(\Omega )\) e ponhamos\(K=\mathrm{supp}(u).\)
	Fixada a família regularizante padrão \((\varphi_{\varepsilon })_{\varepsilon >0}\), se
	\[
		0<\varepsilon <d(K,\mathbb{R}^{n}\setminus \Omega ). \qquad (\#')
	\]
	então
	\[
		\mathrm{supp}(u*\varphi_{\varepsilon })\subseteq \Omega ,
	\]
	e, como \(u*\varphi_{\varepsilon }\in \mathcal{C}^{\infty}(\mathbb{R}^{n})\), segue que
	\[
		u*\varphi_{\varepsilon }\in \mathcal{C}_{c}^{\infty}(\Omega ).
	\]
	Além disso, como \(u\in \mathcal{E}'(\Omega )\) pode ser vista como uma distribuição em \(\mathcal{D}'(\mathbb{R}^{n})\), pelo caso já provado em \(\mathbb{R}^{n}\),
	\[
		u*\varphi_{\varepsilon }\stackrel{\varepsilon \to 0^{+}}\longrightarrow u\quad \text{em }\mathcal{D}'(\mathbb{R}^{n}),
	\]
	e, em particular,
	\[
		u*\varphi_{\varepsilon }\stackrel{\varepsilon \to 0^{+}}\longrightarrow u\quad \text{em }\mathcal{D}'(\Omega ).
	\]
	Isto mostra que \(\mathcal{C}_{c}^{\infty}(\Omega )\) é denso em \(\mathcal{E}'(\Omega )\); combinando com (4), concluímos a densidade em \(\mathcal{D}'(\Omega )\).

	Só falta justificar a afirmação usada em \((\#')\). Por definição,
	\[
		u*\varphi_{\varepsilon }(x)=\langle u, \varphi_{\varepsilon }(x-\cdot )\rangle
		=\langle u(y), \varphi_{\varepsilon }(x-y)\rangle.
	\]
	Como
	\[
		\mathrm{supp}_{y}\bigl(\varphi_{\varepsilon }(x-y)\bigr)
		=x-\mathrm{supp}(\varphi_{\varepsilon })
		=\overline{B}(x;\varepsilon ),
	\]
	se \(\overline{B}(x;\varepsilon )\cap \mathrm{supp}(u)=\emptyset \), então
	\[
		u*\varphi_{\varepsilon }(x)=0.
	\]
	Consequentemente,
	\[
		d(x,K)>\varepsilon \Rightarrow u*\varphi_{\varepsilon }(x)=0,
	\]
	e portanto
	\[
		\mathrm{supp}(u*\varphi_{\varepsilon })\subseteq \overline{O}_{\varepsilon }(K)
		\coloneqq \{x\in \mathbb{R}^{n}:d(x,K)\leq \varepsilon \}.
	\]
	Escolhendo \(\varepsilon \) como em \((\#')\), temos \(\overline{O}_{\varepsilon }(K)\subseteq \Omega \), concluindo a prova. \qedsymbol
\end{proof*}
\end{document}
